\section{Limitaciones de MDLM: el problema del caché KV}

\subsection{Repaso del caché KV}

En los modelos AR, el \textbf{caché KV} (Key-Value Caching) es la técnica clave para acelerar la inferencia.

\begin{knowledgebox}{Principio de funcionamiento del caché KV}
En un Transformer AR:
\begin{enumerate}
    \item Debido al uso de atención causal, los valores de activación de cada token dependen únicamente del contexto izquierdo
    \item Los valores de activación de los tokens ya generados \textbf{nunca cambian}, porque los tokens futuros no los afectan
    \item Por lo tanto, se puede almacenar en caché (congelar) las claves y valores calculados
    \item Al generar un nuevo token, solo es necesario realizar una propagación hacia adelante para un \textbf{único token}, no para todo el contexto
\end{enumerate}
Esto permite que los modelos AR tengan un costo computacional muy bajo por paso, incluso si requieren miles de pasos de decodificación.
\end{knowledgebox}

\subsection{Por qué MDLM no admite caché KV}

MDLM utiliza atención bidireccional, lo que significa que \textbf{los valores de activación de cada token dependen de todos los demás tokens}.

Considere el siguiente escenario:
\begin{enumerate}
    \item Estado inicial: $[\text{M}, \text{M}, \text{M}, \text{M}]$
    \item Después del primer paso de desruido: $[\text{M}, \text{B}, \text{M}, \text{D}]$ (se revelan las posiciones 2 y 4)
    \item Después del segundo paso de desruido: $[A, \text{B}, \text{M}, \text{D}, \text{M}, \text{F}]$ (se revelan las posiciones 1 y 6)
\end{enumerate}

En el segundo paso, aunque $\text{B}$ y $\text{D}$ no han cambiado, debido a la aparición de $A$ y $F$, bajo atención bidireccional los valores de activación de $\text{B}$ y $\text{D}$ cambiarán. Por lo tanto, \textbf{no se puede congelar el valor de activación en ninguna posición}.

\begin{warningbox}{Cuello de botella de inferencia de MDLM}
Aunque MDLM puede generar texto con muchos menos pasos que la longitud del secuencia (por ejemplo, generar 1000 palabras con 400-500 pasos), cada paso requiere realizar una propagación hacia adelante completa para \textbf{todo el secuencia}. Para secuencias largas, la complejidad $O(n^2)$ de la atención hace que el costo por paso sea mucho mayor que el de un solo paso en un modelo AR (con caché KV). Por lo tanto, en términos de tiempo real de inferencia, MDLM podría ser incluso más lento que un modelo AR.
\end{warningbox}

\subsection{Resumen de la sección}

La ausencia del caché KV es la limitación práctica más grande de MDLM. Aunque la atención bidireccional otorga al modelo una capacidad de percepción global, también provoca que los valores de activación se actualicen globalmente ante cambios en las entradas, lo que obstaculiza el uso del caché KV.

\newpage
